\documentclass[12pt,a4paper]{article}
\usepackage[utf8]{vietnam}
\usepackage{amsmath,amssymb}
\usepackage{geometry}
\geometry{
  top=12mm,
  bottom=12mm,
  left=14mm,
  right=14mm,
  footskip=6mm,
  headheight=14pt
}

\usepackage{tikz}
\usetikzlibrary{arrows.meta,calc,shapes.geometric}
\usepackage[most]{tcolorbox}
\usepackage{fontawesome}
\usepackage{fancyhdr}
\usepackage{enumitem}

\pagestyle{fancy}
\fancyhf{}
\fancyhead[L]{\footnotesize\color{teal!80!black}\faBook\ \textbf{SỔ TAY KIẾN THỨC TOÁN 3}}
\fancyhead[R]{\footnotesize\color{teal!80!black}\textbf{Lớp Toán Thầy Lĩnh -- loptoanthaylinh.vn}}
\fancyfoot[C]{\footnotesize\color{gray}-- Sổ tay kiến thức: Bài 23. Nhân số có hai chữ số với số có một chữ số --}
\renewcommand{\headrulewidth}{0.4pt}
\renewcommand{\footrulewidth}{0pt}

\setlength{\parindent}{0pt}
\setlength{\parskip}{0pt}

\begin{document}

% --- BANNER BÀI HỌC ---
\begin{tcolorbox}[
  enhanced,
  colback=teal!8!white,
  colframe=teal!70!black,
  arc=3mm, boxrule=1.2pt,
  left=3mm, right=3mm, top=2mm, bottom=2mm,
  halign=center
]
  {\LARGE\bfseries\color{teal!90!black}\faBookmark\ BÀI 23: NHÂN SỐ CÓ HAI CHỮ SỐ VỚI SỐ CÓ MỘT CHỮ SỐ}\\[1mm]
  {\itshape\footnotesize\color{black!75} Kĩ thuật đặt tính nhân thẳng cột, nhân không nhớ và có nhớ 1 lần trên 1 trang A4}
\end{tcolorbox}

\vspace{-2mm}

\noindent
\begin{minipage}[t]{0.49\textwidth}\vspace{0pt}
\begin{tcolorbox}[
  enhanced,
  colback=blue!4!white,
  colframe=blue!75!black,
  arc=2mm, boxrule=0.8pt,
  left=2mm, right=2mm, top=1.8mm, bottom=1.8mm,
  title={\footnotesize\bfseries\color{white}\faTimes\ 1. Quy trình Đặt tính rồi tính}
]
{\footnotesize
\begin{itemize}\setlength{\itemsep}{1pt}\setlength{\leftskip}{-2.5mm}
  \item \textbf{Đặt tính:} Viết thừa số thứ hai thẳng cột dưới chữ số hàng \textbf{đơn vị} của thừa số thứ nhất. Viết dấu $\times$, kẻ vạch ngang.
  \item \textbf{Thứ tự nhân:} Lấy thừa số thứ hai nhân lần lượt với từng chữ số của thừa số thứ nhất từ \textbf{phải sang trái}:
  \begin{enumerate}\setlength{\itemsep}{0.5pt}\setlength{\leftskip}{-2.5mm}
    \item Nhân hàng đơn vị trước.
    \item Nhân hàng chục sau.
  \end{enumerate}
\end{itemize}
}
\end{tcolorbox}
\end{minipage}
\hfill
\begin{minipage}[t]{0.49\textwidth}\vspace{0pt}
\begin{tcolorbox}[
  enhanced,
  colback=teal!4!white,
  colframe=teal!75!black,
  arc=2mm, boxrule=0.8pt,
  left=2mm, right=2mm, top=1.8mm, bottom=1.8mm,
  title={\footnotesize\bfseries\color{white}\faMagic\ 2. Phép nhân có nhớ}
]
{\footnotesize
\begin{itemize}\setlength{\itemsep}{1pt}\setlength{\leftskip}{-2.5mm}
  \item Khi nhân hàng đơn vị có kết quả $\ge 10$:
  \begin{itemize}\setlength{\itemsep}{0.5pt}\setlength{\leftskip}{-2.5mm}
    \item Viết chữ số hàng đơn vị vào kết quả.
    \item \textbf{Nhớ} số chục sang hàng chục.
  \end{itemize}
  \item Khi nhân hàng chục: Lấy thừa số nhân với hàng chục, rồi \textbf{cộng thêm số nhớ} vào kết quả:
  $$\text{Tích hàng chục} = (\text{Thừa số} \times \text{Hàng chục}) + \text{Số nhớ}$$
\end{itemize}
}
\end{tcolorbox}
\end{minipage}

\vspace{-2.5mm}
% --- KHỐI 3: SƠ ĐỒ TƯ DUY TIKZ ---
\noindent
\begin{tcolorbox}[
  enhanced,
  colback=orange!4!white,
  colframe=orange!80!black,
  arc=2mm, boxrule=0.8pt,
  left=2mm, right=2mm, top=1.5mm, bottom=1.5mm,
  title={\footnotesize\bfseries\color{white}\faSitemap\ SƠ ĐỒ TƯ DUY: TIẾN TRÌNH 2 BƯỚC NHÂN TỪ PHẢI SANG TRÁI}
]
\centering
\begin{tikzpicture}[
  box/.style={draw, rounded corners=2pt, align=center, font=\scriptsize\bfseries, inner sep=3pt, minimum height=0.85cm},
  arr/.style={-{Stealth[length=2.5mm]}, thick}
]
  \node[box, fill=teal!15, draw=teal!80!black, text width=3.3cm] (n1) at (0, 0) {BƯỚC 1: NHÂN ĐƠN VỊ\\$b \times c$\\Viết đơn vị, nhớ chục};
  \node[box, fill=blue!12, draw=blue!80!black, text width=3.3cm] (n2) at (4.7, 0) {BƯỚC 2: NHÂN HÀNG CHỤC\\$a \times c$\\Cộng thêm số nhớ};
  \node[box, fill=purple!12, draw=purple!80!black, text width=3.3cm] (n3) at (9.4, 0) {BƯỚC 3: KẾT QUẢ\\Ghi tích hoàn chỉnh\\Kiểm tra tính hợp lí};

  \draw[arr, teal!80!black] (n1) -- (n2);
  \draw[arr, blue!80!black] (n2) -- (n3);
\end{tikzpicture}
\end{tcolorbox}

\vspace{-2.5mm}
% --- KHỐI 4: VÍ DỤ & MẸO TRÁNH BẪY ---
\noindent
\begin{minipage}[t]{0.49\textwidth}\vspace{0pt}
\begin{tcolorbox}[
  enhanced,
  colback=green!4!white,
  colframe=green!60!black,
  arc=2mm, boxrule=0.8pt,
  left=2mm, right=2mm, top=1.8mm, bottom=1.8mm,
  title={\footnotesize\bfseries\color{white}\faCheckCircle\ Ví dụ thực hành chuẩn mực}
]
{\footnotesize
\textbf{1. Phép nhân không nhớ:}\\
\begin{tabular}{r@{\quad}l}
  $23$ & $\bullet\ 3 \times 3 = 9$, viết $9$.\\
$\times\phantom{0}3$ & $\bullet\ 3 \times 2 = 6$, viết $6$.\\
\cline{1-1}
  $\mathbf{69}$ & (Kết quả: $69$).
\end{tabular}\\[1mm]
\textbf{2. Phép nhân có nhớ:}\\
\begin{tabular}{r@{\quad}l}
  $26$ & $\bullet\ 3 \times 6 = 18$, viết $8$ nhớ $1$.\\
$\times\phantom{0}3$ & $\bullet\ 3 \times 2 = 6$, thêm $1$ là $7$, viết $7$.\\
\cline{1-1}
  $\mathbf{78}$ & (Kết quả: $78$).
\end{tabular}
}
\end{tcolorbox}
\end{minipage}
\hfill
\begin{minipage}[t]{0.49\textwidth}\vspace{0pt}
\begin{tcolorbox}[
  enhanced,
  colback=red!4!white,
  colframe=red!70!black,
  arc=2mm, boxrule=0.8pt,
  left=2mm, right=2mm, top=1.8mm, bottom=1.8mm,
  title={\footnotesize\bfseries\color{white}\faExclamationTriangle\ Mẹo tránh bẫy trắc nghiệm}
]
{\footnotesize
\begin{itemize}\setlength{\itemsep}{1.5pt}\setlength{\leftskip}{-2.5mm}
  \item \textbf{Bẫy quên cộng số nhớ:}\\
  Ví dụ $26 \times 3$: Rất nhiều bạn tính $3 \times 6 = 18$ nhớ $1$, nhưng đến hàng chục lấy $3 \times 2 = 6$ rồi ghi luôn $68$ (quên cộng $1$ nhớ)!
  \item \textbf{Bẫy cộng nhớ trước khi nhân:}\\
  Phải lấy $3 \times 2 = 6$ rồi mới cộng $1$ thành $7$. Tuyệt đối không được lấy hàng chục $2$ cộng $1$ nhớ thành $3$ rồi mới nhân $3 \times 3 = 9$ (SAI NẶNG!).
\end{itemize}
}
\end{tcolorbox}
\end{minipage}


\end{document}
